\section{Shift-averaged discrepancy and a prefix criterion}\label{sec:framework}
Let $K$ be a real symmetric reproducing kernel on $[0,1)^d$ with Hilbert
space $\Hspace$. Assume that point evaluation and integration are bounded,
and that the integrals below exist and can be interchanged with the finite
sums. For deterministic nodes $P_n=(\vx_0,\ldots,\vx_{n-1})$, define
\[
 D_K^2(P_n)=\sup_{\|f\|_{\Hspace}\le1}|I(f)-Q_n(f)|^2.
\]
The reproducing property gives the familiar kernel expression
\begin{align}\label{eq:rkhs-discrepancy}
 D_K^2(P_n)
 &=\int\!\!\int K(\vx,\boldsymbol y)\dif\vx\dif\boldsymbol y
 -\frac2n\sum_{i=0}^{n-1}\int K(\vx_i,\boldsymbol y)\dif\boldsymbol y\notag\\
 &\quad+\frac1{n^2}\sum_{i,j=0}^{n-1}K(\vx_i,\vx_j).
\end{align}
All unmarked integrals in this section are over the unit cube of the
appropriate dimension.

For $\vDelta\sim\mathrm{Unif}[0,1)^d$, independent of the nodes, put
\[
 e_n^2(P_n)=\E_{\vDelta}\left[D_K^2\bigl(
 (\fracpart{\vx_i+\vDelta})_{i=0}^{n-1}\bigr)\right],
 \quad
 \kt(\boldsymbol t)=\int K(\fracpart{\boldsymbol y+\boldsymbol t},
 \boldsymbol y)\dif\boldsymbol y,
 \quad \overline K=\int\kt(\boldsymbol t)\dif\boldsymbol t.
\]
Uniform translation modulo one yields
\begin{equation}\label{eq:shift-esd}
 e_n^2(P_n)=-\overline K+
 \frac1{n^2}\sum_{i,j=0}^{n-1}
 \kt(\fracpart{\vx_i-\vx_j}).
\end{equation}
Indeed, each shifted node is uniform, so both the constant integral term
and each averaged single-node integral in~\eqref{eq:rkhs-discrepancy}
equal $\overline K$. The averaged pair term is the translation-averaged
kernel at the unshifted difference. In particular, a random shift belongs
in the integration rule, but must not be inserted again into the differences
on the right side of~\eqref{eq:shift-esd}.

This is an expectation of a \emph{squared worst-case error}. It is not a
confidence interval, nor does a small value establish a finite-sample
probabilistic error guarantee for an arbitrary integrand outside the space.
For $f\in\Hspace$, the deterministic worst-case inequality does imply
$\E_{\vDelta}|I(f)-Q_n(f)|^2\le e_n^2(P_n)\|f\|_{\Hspace}^2$.

\subsection{Product kernels and two kinds of weights}
We specialize to a periodic product kernel
\begin{equation}\label{eq:kernel}
 \kt_d(\boldsymbol t)=\prod_{j=1}^d
 \bigl[1+\gamma_j B_2(\fracpart{t_j})\bigr],
 \qquad B_2(t)=t^2-t+\tfrac16,
 \qquad \gamma_j>0.
\end{equation}
Here $\overline K=1$, since $B_2$ has zero mean. The Fourier coefficients
of $B_2(\fracpart{t})$ are positive for nonzero integer frequencies, so
this defines a positive-definite periodic kernel. Its normalization is the
one used in the talk's computations. Rescaling $B_2$, for example by
$2\pi^2$, also rescales the effective coordinate weights and cannot be done
silently when comparing figures or generators.

For sample-size weights $w_1,\ldots,w_N\ge0$, define
\begin{equation}\label{eq:wesd}
 \W_N(P)=\sum_{n=1}^N w_n e_n^2(P_n),
 \qquad\text{with }w_n=n\text{ in the constructions below}.
\end{equation}
The same underlying sequence supplies every prefix $P_n$. Choosing the
weights is part of specifying the objective; small WSESD is an aggregate
property and does not mean that every individual prefix beats every
comparator. For $w_n>0$ it gives the elementary bound
$e_n^2(P_n)\le\W_N(P)/w_n$, which by itself may be quite conservative.

One interpretation makes the role of these weights explicit. If a random
stopping size $T$ has probabilities $\Pr(T=n)=p_n$, then
$\E[e_T^2]=\sum_n p_n e_n^2$. Alternatively,
$\E[T^2e_T^2]=\sum_n p_n n^2 e_n^2$ emphasizes squared error on the
scale $TQ_T$. With $p_n=(nH_N)^{-1}$ and
$H_N=\sum_{n=1}^N n^{-1}$, the latter is $\W_N/H_N$ for $w_n=n$.
This is an interpretation of the objective, not an assertion that the
applications in the talk follow this stopping distribution.

\par\smallskip
\noindent\codex{[Framework verification: Identify the intended periodic
function space and a supporting reference for this kernel normalization.
Explain the choices of sample-size weights and coordinate weights used in the
final study; the stopping-distribution interpretation above is a derived
interpretation, not evidence about an application.]}
\par\smallskip
